\chapter{Variables de prix et de volume}\label{ch:rs:price-and-volume-features}

Le facteur momentum publié dans la bibliothèque de données de Kenneth French, comme la plupart des signaux de momentum, exclut le mois le plus récent~: il
classe les actions selon leur rendement entre il y a douze mois et il y a un mois. La raison tient à un effet, non à une convention. Les rendements mensuels des actions présentent une autocorrélation
d'ordre un fortement négative (Jegadeesh, 1990), si bien que les gagnants du dernier mois en rendent une partie, et inclure ce
mois mêle un renversement à un signal de continuation. Sur le marché synthétique du livre, l'exclusion produit l'effet inverse~: le momentum sur
les douze mois complets a un \omterm{def:rs:anatomy-of-a-predictor:ic}{coefficient d'information} à un mois de 0.034, la version avec exclusion 0.030. Le renversement du simulateur
dure un jour, non un mois, et sur un marché sans renversement mensuel l'exclusion jette de l'information. Les détails
d'une variable encodent des affirmations sur le marché~: ce chapitre construit la famille standard des variables de prix et de volume, mesure chacune
sur le marché synthétique, et indique quelle affirmation chacune formule.

\section{Rendements à plusieurs échelles}

\begin{definition}[Variable de rendement, période d'exclusion]\label{def:rs:price-and-volume-features:return}
Une \emph{variable de rendement}\index{variable de rendement} est le rendement composé d'un titre sur une fenêtre qui se termine au plus tard
à l'instant de décision, utilisée comme prédicteur. Sa \emph{période d'exclusion}\index{période d'exclusion} est l'écart entre la fin de la fenêtre et l'instant de décision~:
une variable de momentum 12-1 utilise le rendement entre 252 et 21 jours de bourse plus tôt, en excluant les 21 derniers.
\end{definition}

Les \omterm{def:rs:price-and-volume-features:return}{variables de rendement} comptent parmi les plus anciens prédicteurs et restent les plus utilisées~: une fenêtre courte (d'un jour à un mois) pour le renversement, une fenêtre longue
(de trois à douze mois) pour le momentum, de trois à cinq ans pour le renversement de long terme (De Bondt et Thaler, 1985). Jegadeesh et Titman
(1993) ont montré qu'acheter les gagnants passés et vendre les perdants passés sur trois à douze mois rapportait des rendements significatifs~; le livre~8
étudie ces stratégies. Ici, la question est celle de la mesure.

Le tableau mesure les douze variables du chapitre sur le marché synthétique~: l'IC de rang moyen par rapport au jour, à la semaine et au mois suivants,
chacune échantillonnée une fois par horizon, afin que les cibles ne se chevauchent pas et que les $t$-statistiques naïves soient valides. Le simulateur introduit un
renversement d'un jour et une dérive persistante (chapitre~5), et le tableau retrouve les deux~: les variables de renversement obtiennent un score aux horizons courts, le momentum
et les variables de tendance aux horizons longs. Il n'introduit rien sur la volatilité, le volume ou la liquidité, et ces variables
ne se distinguent pas de zéro. Ce n'est pas un constat sur les marchés réels, où chacun a sa littérature~; c'est un contrôle que la
mesure n'invente pas d'effets.

\begin{center}
\small
\begin{tabular}{lrrr}
\toprule
variable (connue à la clôture de $t$) & 1 jour & 5 jours & 21 jours \\
\midrule
renversement, 1 jour & 0.0378 (16.2) & 0.0150 (3.0) & $-0.0058$ ($-0.6$) \\
renversement, 5 jours & 0.0176 (8.1) & 0.0085 (1.8) & $-0.0162$ ($-1.8$) \\
renversement, 21 jours & 0.0022 (1.0) & $-0.0117$ ($-2.6$) & $-0.0285$ ($-2.9$) \\
momentum 12-1 & 0.0076 (1.4) & 0.0167 (1.4) & 0.0296 (1.2) \\
momentum 12 & 0.0059 (1.1) & 0.0177 (1.5) & 0.0340 (1.4) \\
momentum 6-1 & 0.0067 (1.8) & 0.0140 (1.7) & 0.0285 (1.6) \\
faible volatilité (21 jours) & $-0.0034$ ($-2.5$) & $-0.0033$ ($-1.1$) & $-0.0023$ ($-0.5$) \\
\omterm{def:rs:price-and-volume-features:volume}{surprise de volume} & 0.0009 (1.3) & $-0.0007$ ($-0.5$) & $-0.0008$ ($-0.3$) \\
\omterm{def:rs:price-and-volume-features:volume}{illiquidité d'Amihud} & $-0.0004$ ($-0.3$) & $-0.0023$ ($-0.7$) & $-0.0055$ ($-0.8$) \\
faible \omterm{def:rs:price-and-volume-features:volume}{taux de rotation des titres} & 0.0011 (1.5) & 0.0026 (1.7) & 0.0049 (1.7) \\
croisement 5/20 & 0.0027 (1.2) & 0.0133 (2.9) & 0.0218 (2.4) \\
croisement 50/200 & 0.0084 (2.0) & 0.0173 (1.9) & 0.0324 (1.7) \\
\bottomrule
\end{tabular}
\end{center}
\noindent{\small IC de rang moyen (et $t$-statistique) de chaque variable par rapport aux rendements des 1, 5 et 21 jours suivants sur le marché synthétique,
dix ans, une observation par horizon. Données~: \texttt{firm.synthmkt}, graine 1.}

\begin{proposition}[Ce que vaut l'exclusion]\label{prop:rs:price-and-volume-features:skip}
On écrit le log-rendement sur douze mois comme somme de la partie 12-1 et de celle du dernier mois. Si les deux sont standardisées de la même manière, l'IC
du rendement complet est, au premier ordre, une moyenne des IC des deux parties pondérée par leurs variances. L'exclusion relève l'IC exactement lorsque le
rendement du dernier mois a un IC plus faible que la partie 12-1, en particulier lorsqu'il est négatif (un renversement mensuel).
\end{proposition}

\begin{proof}
Pour une somme $s = a + b$ et une cible $y$, $\operatorname{corr}(s, y) = (\sigma_a\operatorname{corr}(a, y) + \sigma_b\operatorname{corr}(b,
y))/\sigma_s$, avec $\sigma_s \le \sigma_a + \sigma_b$~; le signal complet bat $a$ seul si et seulement si $\operatorname{corr}(b, y)$ est
suffisamment grand par rapport à $\operatorname{corr}(a, y)$, et il perd dès que $\operatorname{corr}(b, y) < 0$.
\end{proof}

Sur le marché synthétique, le rendement propre du dernier mois a un IC de $+0.0285$ avec le mois suivant (la variable de renversement
donne $-0.0285$ au signe près)~: il y a continuation, si bien que l'année complète l'emporte. Dans les données américaines de Jegadeesh (1990), le dernier mois s'inverse, et l'exclusion l'emporte.
La \Cref{fig:rs:price-and-volume-features:horizon} montre l'IC de chaque variable en fonction de la longueur de sa cible~: le renversement d'un jour est
épuisé en moins d'une semaine, le momentum continue de s'accumuler pendant trois mois.

\begin{omfigure}
\begin{tikzpicture}
\begin{semilogxaxis}[width=10.5cm,height=5.2cm,xlabel={horizon de la cible $h$ (jours de bourse, échelle logarithmique)},ylabel={IC de rang moyen},
  tick label style={font=\small},label style={font=\small},xmin=0.9,xmax=70,ymin=-0.03,ymax=0.045,grid=major,grid style={gray!25},
  xtick={1,2,5,10,21,42,63},xticklabels={1,2,5,10,21,42,63},scaled y ticks=false,
  yticklabel style={/pgf/number format/fixed,/pgf/number format/precision=2},legend columns=2,legend cell align=left,
  legend style={font=\footnotesize,at={(0.5,-0.32)},anchor=north,draw=none,fill=none,
  /tikz/every even column/.append style={column sep=8pt}},table/col sep=comma]
\addplot[omThm,very thick,mark=*] table[x=h,y=rev1]{figdata/research/07-price-and-volume-features/horizon.csv};
\addplot[omThm!55,thick,dashed,mark=o] table[x=h,y=rev21]{figdata/research/07-price-and-volume-features/horizon.csv};
\addplot[omDef,very thick,mark=square*] table[x=h,y=mom121]{figdata/research/07-price-and-volume-features/horizon.csv};
\addplot[omDef!55,thick,dashed,mark=square] table[x=h,y=mom12]{figdata/research/07-price-and-volume-features/horizon.csv};
\addplot[black!60,thin,sharp plot] coordinates {(0.9,0) (70,0)};
\legend{renversement 1 jour,renversement 21 jours,momentum 12-1,momentum 12}
\end{semilogxaxis}
\end{tikzpicture}

\omcaption{IC de rang moyen de quatre \omterm{def:rs:price-and-volume-features:return}{variables de rendement} par rapport aux rendements des $h$ jours suivants sur le marché synthétique, évalué chaque jour.
Le renversement d'un jour atteint son maximum à un jour et disparaît après deux semaines~; le renversement du dernier mois est négatif à chaque horizon (le mois
continue)~; le momentum atteint environ 0.04 à un trimestre, un peu plus haut sans l'exclusion. Données~: \texttt{firm.synthmkt}, graine 1.}\label{fig:rs:price-and-volume-features:horizon}
\end{omfigure}

\section{Volatilité et étendue}

\begin{definition}[Estimateur de volatilité par l'étendue]\label{def:rs:price-and-volume-features:range}
Un \emph{estimateur de volatilité par l'étendue}\index{estimateur de volatilité par l'étendue} estime la variance du rendement d'un jour à partir de son ouverture,
son plus haut, son plus bas et son cours de clôture, plutôt que des seuls rendements clôture à clôture.
\end{definition}

\begin{definition}[Estimateurs de Parkinson, de Garman--Klass et de Yang--Zhang]\label{def:rs:price-and-volume-features:estimators}
Avec $H, L, O, C$ le plus haut, le plus bas, l'ouverture et la clôture d'une journée, l'\emph{estimateur de Parkinson}\index{estimateur de Parkinson} de la variance journalière est
$(\ln H/L)^2/(4\ln 2)$~; l'\emph{estimateur de Garman--Klass}\index{estimateur de Garman--Klass} est $\tfrac12(\ln H/L)^2 - (2\ln 2 -
1)(\ln C/O)^2$~; celui de Rogers--Satchell est $\ln(H/C)\ln(H/O) + \ln(L/C)\ln(L/O)$. L'\emph{estimateur de
Yang--Zhang}\index{estimateur de Yang--Zhang} sur $n$ jours ajoute la variance empirique des rendements de nuit, $k$ fois celle des
rendements ouverture-clôture et $1 - k$ fois le terme moyen de Rogers--Satchell, avec $k = 0.34/(1.34 + (n+1)/(n-1))$.
\end{definition}

Pour un mouvement brownien sans dérive observé en continu, $\E[(\ln H/L)^2] = 4\ln2\,\sigma^2$, ce qui est la normalisation de Parkinson
(Parkinson, 1980)~; l'étendue porte beaucoup plus d'information sur $\sigma$ que la seule clôture. Chaque estimateur postérieur corrige une
faiblesse du précédent. Garman et Klass (1980) combinent l'étendue avec le rendement ouverture-clôture~; Rogers et Satchell (1991) sont
sans biais quelle que soit la dérive~; Yang et Zhang (2000) ajoutent l'écart de nuit, que nulle étendue intrajournalière ne peut voir.

La simulation du chapitre sur 4\,000 jours de trajectoires browniennes à la minute les évalue (\cref{fig:rs:price-and-volume-features:race}).
Sans dérive ni écart, la variance de l'\omterm{def:rs:price-and-volume-features:estimators}{estimateur de Parkinson} est 5.2 fois plus petite que celle du carré du rendement clôture-à-clôture,
celle de Garman--Klass 7.6 fois, celle de Rogers--Satchell 5.7 fois~; l'\omterm{def:rs:price-and-volume-features:estimators}{estimateur de Yang--Zhang} sur vingt jours est 7.5 fois plus efficace que la
variance clôture-à-clôture sur vingt jours. Chaque estimateur par l'étendue sous-estime d'environ 8 à 10~\%, car une étendue observée en 390 points est plus étroite
que l'étendue continue~: la surveillance discrète introduit un biais que les formules ignorent. Une dérive égale à un demi-écart-type par jour
gonfle le carré du rendement clôture-à-clôture de 25~\% et celui de Parkinson de 2~\%, et laisse Rogers--Satchell et Yang--Zhang inchangés~; un écart de nuit
portant 30~\% de la variance fait que les trois estimateurs par l'étendue d'une seule journée ne voient que 63 à 64~\% de celle-ci, et seul Yang--Zhang retrouve
la totalité.

\begin{omfigure}
\begin{tikzpicture}
\begin{axis}[width=11cm,height=5.4cm,ybar,bar width=8pt,ylabel={estimation moyenne / variance vraie},tick label style={font=\small},
  label style={font=\small},xmin=-0.6,xmax=4.6,ymin=0,ymax=1.35,xtick={0,1,2,3,4},
  xticklabels={{clôture-à-\\clôture},{Parkinson},{Garman--\\Klass},{Rogers--\\Satchell},{Yang--\\Zhang}},
  xticklabel style={align=center,font=\footnotesize},
  grid=major,grid style={gray!25},legend columns=3,legend style={font=\footnotesize,at={(0.5,-0.36)},anchor=north,draw=none,
  fill=none,/tikz/every even column/.append style={column sep=8pt}},table/col sep=comma]
\addplot[fill=omDef!55,draw=omDef] table[x=k,y=base]{figdata/research/07-price-and-volume-features/race.csv};
\addplot[fill=omThm!45,draw=omThm] table[x=k,y=drift]{figdata/research/07-price-and-volume-features/race.csv};
\addplot[fill=black!25,draw=black!60] table[x=k,y=gaps]{figdata/research/07-price-and-volume-features/race.csv};
\addplot[black!70,thin,sharp plot,forget plot] coordinates {(-0.6,1) (4.6,1)};
\legend{sans dérive ni écart,dérive égale à un demi-écart-type journalier,écart de nuit~: 30~\% de la variance}
\end{axis}
\end{tikzpicture}

\omcaption{Moyenne de chaque estimateur de variance journalière relative à la variance de diffusion vraie, sur 4\,000 jours simulés de 390 pas
d'une minute. Tous les estimateurs par l'étendue sous-estiment de 8 à 10~\% à cause de la surveillance discrète~; une dérive gonfle la clôture-à-clôture et Parkinson~; un écart de nuit
est invisible aux estimateurs par l'étendue d'une seule journée et retrouvé par Yang--Zhang. Données~: le module du chapitre, avec graine fixée.}\label{fig:rs:price-and-volume-features:race}
\end{omfigure}

\section{Volume et liquidité}

\begin{definition}[Surprise de volume, illiquidité d'Amihud, taux de rotation des titres]\label{def:rs:price-and-volume-features:volume}
La \emph{surprise de volume}\index{surprise de volume} d'un titre sur un jour donné est le logarithme de son volume divisé par son volume moyen
sur une fenêtre précédente. L'\emph{illiquidité d'Amihud}\index{illiquidité d'Amihud} est la moyenne sur une fenêtre de la valeur absolue du rendement journalier
divisée par la valeur échangée du jour~: le mouvement de prix par unité de monnaie échangée (Amihud, 2002). Le \emph{taux de rotation des titres}%
\index{taux de rotation des titres} est le volume échangé divisé par le nombre d'actions en circulation, moyenné sur une fenêtre.
\end{definition}

Les variables de volume mesurent l'attention et la liquidité, et chacune s'appuie sur une thèse. Gervais, Kaniel et Mingelgrin (2001) ont constaté que
les actions dont le volume est anormalement élevé sur un jour ou une semaine ont tendance à monter le mois suivant (la prime de rendement des volumes élevés)~: le volume
attire l'attention, et l'attention des investisseurs fait bouger les prix. La mesure d'Amihud est un proxy bon marché de l'impact sur les prix, et les actions peu liquides
devraient rémunérer ce risque par une prime. Aucun de ces mécanismes n'est présent dans le simulateur du livre, dont le volume suit l'ampleur de chaque
jour et l'état de la volatilité, mais ne porte aucune information sur les rendements futurs~; le tableau le montre précisément. Un chercheur qui
veut les étudier a besoin de données réelles, ou d'un simulateur qui les intègre.

\section{Familles techniques comme filtres linéaires}

\begin{definition}[Croisement de moyennes mobiles, filtre linéaire]\label{def:rs:price-and-volume-features:filter}
Un \emph{croisement de moyennes mobiles}\index{croisement de moyennes mobiles} de fenêtres $n < m$ est le logarithme du rapport entre la moyenne
des $n$ derniers prix et la moyenne des $m$ derniers prix~; la règle classique est d'acheter lorsqu'il est positif. Un \emph{filtre
linéaire}\index{filtre linéaire} sur les rendements est un prédicteur de la forme $\sum_{j \ge 0} w_jr_{t-j}$ à poids fixes $w_j$.
\end{definition}

\begin{proposition}[Un croisement est un filtre sur les rendements passés]\label{prop:rs:price-and-volume-features:crossover}
Au premier ordre en les rendements, $\ln(\overline P_n/\overline P_m) = \sum_{j=0}^{m-2} w_jr_{t-j}$ avec $w_j = \frac{(m - 1 - j)}{m} -
\frac{(n - 1 - j)^+}{n}$, où $r_{t-j}$ est le log-rendement d'il y a $j$ jours~: les poids croissent linéairement sur la fenêtre courte, puis décroissent
linéairement jusqu'à zéro à la longueur de la fenêtre longue, et leur somme vaut $\frac{m - n}{2}$.
\end{proposition}

\begin{proof}
Le logarithme d'une moyenne de prix vaut, au premier ordre, la moyenne des logarithmes des prix. La moyenne des $n$ derniers logarithmes de prix vaut $\ln P_t
- \sum_{j<n-1}\frac{n-1-j}{n}r_{t-j}$, puisque le prix il y a $k$ jours vaut $\ln P_t$ moins les $k$ rendements les plus récents. On soustrait la
même expression avec $m$. Les poids somment à $\frac{m-1}{2} - \frac{n-1}{2}$.
\end{proof}

Un croisement est donc un signal de momentum à fenêtre triangulaire qui ignore la plupart des $n$ derniers jours
(\cref{fig:rs:price-and-volume-features:filter}), et sa diversité apparente n'est qu'une diversité de poids sur les rendements passés. Brock, Lakonishok
et LeBaron (1992) ont testé des règles de moyennes mobiles et de fourchettes de cours sur l'indice Dow Jones de 1897 à 1986 et ont trouvé un appui fort pour
elles~; la vision par filtres dit de quoi parle un tel résultat~: l'autocorrélation des rendements de l'indice aux retards que les poids mettent en avant. Sur le
marché synthétique, le croisement 5/20 a un IC de 0.022 à un mois ($t = 2.4$), version lissée du momentum que le simulateur intègre.

\begin{omfigure}
\begin{tikzpicture}
\begin{axis}[width=10.5cm,height=4.8cm,xlabel={retard $j$ du rendement journalier (jours)},ylabel={poids $w_j$},tick label style={font=\small},
  label style={font=\small},xmin=0,xmax=50,ymin=0,ymax=0.9,grid=major,grid style={gray!25},legend columns=2,
  legend style={font=\footnotesize,at={(0.5,-0.32)},anchor=north,draw=none,fill=none,
  /tikz/every even column/.append style={column sep=8pt}},table/col sep=comma]
\addplot[omDef,very thick,const plot mark mid] table[x=lag,y=w520]{figdata/research/07-price-and-volume-features/filter.csv};
\addplot[omThm,very thick,dashed,const plot mark mid] table[x=lag,y=w1050]{figdata/research/07-price-and-volume-features/filter.csv};
\legend{croisement 5/20,croisement 10/50}
\end{axis}
\end{tikzpicture}

\omcaption{Poids qu'un \omterm{def:rs:price-and-volume-features:filter}{croisement de moyennes mobiles} accorde à chaque rendement journalier passé (\cref{prop:rs:price-and-volume-features:crossover})~:
croissants sur la fenêtre courte, puis décroissants linéairement jusqu'à zéro à la fenêtre longue. Le croisement 5/20 somme à 7.5 et culmine au retard 5~; le
10/50 à 20 et culmine au retard 10.}\label{fig:rs:price-and-volume-features:filter}
\end{omfigure}

\section{Fiches de prédicteur}

\begin{predictorcard}{Momentum 12-1}\label{pred:rs:price-and-volume-features:momentum}
\sfield{Définition} $s_{i,t} = \prod_{k=t-251}^{t-21}(1 + r_{i,k}) - 1$~: le rendement entre 252 et 21 jours de bourse en arrière, classé.
\sfield{Données et horodatage} Clôtures journalières jusqu'à 21 jours avant la date de décision~; \omterm{def:rs:universe-symbology-and-corporate-actions:delisting}{rendements de radiation} inclus.
\sfield{Justification} Sous-réaction à l'information et diffusion lente~; le dernier mois est exclu car il s'inverse dans les données américaines
(Jegadeesh, 1990).
\sfield{Horizon et demi-vie} Un à trois mois. Sur le marché synthétique, l'IC avec le rendement d'un seul jour vaut 0.0076 le lendemain et
environ 0.004 six mois plus tard~; un ajustement exponentiel donne une demi-vie de 260 jours.
\sfield{Normalisation} Rang en coupe transversale~; neutralisé par secteur lorsque le momentum sectoriel n'est pas souhaité.
\sfield{Modes de défaillance} Les krachs du momentum surviennent après les rebonds de marché (livre~8)~; l'encombrement (crowding)~; les coûts de rotation.
\sfield{Sources} Jegadeesh et Titman (1993)~; \texttt{rs\_features.lag\_profile} et \texttt{half\_life}.
\end{predictorcard}

\begin{predictorcard}{Surprise de volume}\label{pred:rs:price-and-volume-features:volume}
\sfield{Définition} $s_{i,t} = \ln\bigl(V_{i,t}/\tfrac1{21}\sum_{k=1}^{21}V_{i,t-k}\bigr)$, classé.
\sfield{Données et horodatage} Volumes journaliers jusqu'à la clôture de $t$ (le volume consolidé n'est connu qu'après la clôture).
\sfield{Justification} Un volume inhabituel attire l'attention et les acheteurs~; la prime de rendement des volumes élevés (Gervais, Kaniel et Mingelgrin, 2001).
\sfield{Horizon et demi-vie} Un mois dans la littérature. Non mesurable sur le marché synthétique, qui n'intègre aucun effet de volume~: son IC
vaut 0.0009 ($t = 1.3$) à un jour et zéro ensuite.
\sfield{Normalisation} Rang~; la moyenne exclut le jour courant.
\sfield{Modes de défaillance} Le volume lié aux nouvelles porte la direction de la nouvelle, non l'attention~; les jours de rééquilibrage d'indice et d'échéance d'options gonflent le volume
mécaniquement.
\sfield{Sources} Gervais, Kaniel et Mingelgrin (2001)~; le tableau du chapitre.
\end{predictorcard}

\section{Tutoriel~: la bibliothèque de variables}\label{tut:rs:price-and-volume-features:library}

\begin{tutorial}
\textbf{Objectif.} Calculer les douze variables du chapitre sur le marché synthétique, démontrer qu'aucune ne lit le futur, et remplir le tableau des IC.
\textbf{État final~:} le tableau~; les \cref{fig:rs:price-and-volume-features:horizon,fig:rs:price-and-volume-features:race,fig:rs:price-and-volume-features:filter}.
\begin{tutsteps}
\item \textbf{Fenêtres sans anticipation.} Toute statistique glissante est construite à partir de sommes cumulées~: la ligne $t$ utilise les lignes $t - w + 1$ à $t$,
et une fenêtre contenant une valeur manquante est manquante.
\omcode{code/firm/features/firm_features.py}{29}{53}{Sommes glissantes, décalages et rendements passés.}\label{lst:rs:price-and-volume-features:rolling}
\item \textbf{Le test de fuite.} Perturber chaque entrée après un instant de coupure et exiger que la variable jusqu'à cet instant reste inchangée~; les
tests d'acceptation l'appliquent à chaque variable, et à une variable qui regarde vers l'avenir, qu'il doit détecter.
\omcode{code/firm/features/firm_features.py}{137}{148}{Le test de fuite.}\label{lst:rs:price-and-volume-features:leak}
\item \textbf{Exécuter} \texttt{ic\_table()}, \texttt{ic\_by\_horizon} pour les quatre \omterm{def:rs:price-and-volume-features:return}{variables de rendement}, \texttt{range\_race} pour les trois cas
et \texttt{fig\_features.py}.
\end{tutsteps}
\textbf{Pour aller plus loin.} Ajouter le renversement à long terme (le rendement entre il y a cinq ans et il y a un an) et vérifier son IC~; remplacer les trajectoires
simulées à la minute par une journée de firm.tape et comparer les \omterm{def:rs:price-and-volume-features:range}{estimateurs de volatilité par l'étendue} sur un prix affecté d'un rebond bid-ask.
\end{tutorial}

\section{Construction~: la bibliothèque de variables}\label{bld:rs:price-and-volume-features:features}

\begin{build}
\bfield{Rôle} Les variables standard sur barres de la petite entreprise, calculées de la même façon partout, avec un test qu'aucune d'elles ne regarde vers l'avenir.
\bfield{Interface} \texttt{past\_return(ret, window, skip)}, \texttt{rolling\_vol}, \texttt{ewma\_vol}, \texttt{parkinson},
\texttt{garman\_klass}, \texttt{rogers\_satchell}, \texttt{yang\_zhang}, \texttt{volume\_surprise}, \texttt{amihud},
\texttt{turnover}, \texttt{ma\_crossover}, \texttt{crossover\_weights}, \texttt{leakage\_test(feature, inputs, cut)}.
\bfield{Règles} Panels (dates, noms)~; la ligne $t$ d'une sortie utilise les lignes d'entrée jusqu'à $t$ seulement~; les fenêtres contenant une valeur manquante sont manquantes~;
les moyennes de volume excluent le jour courant.
\bfield{Tests d'acceptation} \texttt{code/firm/features/tests/}~: rendements passés avec une \omterm{def:rs:price-and-volume-features:return}{période d'exclusion}~; chaque variable passe le test de fuite et une
variable qui regarde vers l'avenir échoue~; les estimateurs par l'étendue à 10~\% près de la variance vraie d'une journée brownienne simulée~; le croisement égale son
filtre au premier ordre.
\bfield{Pour aller plus loin} Variables intrajournalières à partir des barres de firm.tape~; variables sur barres à horloge de volume (chapitre~2)~; un registre qui enregistre la
fenêtre et le \omterm{def:rs:point-in-time-data-and-the-biases:bitemporal}{temps de connaissance} de chaque variable pour la \omterm{def:rs:anatomy-of-a-predictor:card}{fiche de prédicteur}.
\end{build}

\begin{omsources}
\item N.~Jegadeesh, ``Evidence of predictable behavior of security returns'', \emph{Journal of Finance} 45(3), 1990~; N.~Jegadeesh et
S.~Titman, ``Returns to buying winners and selling losers'', \emph{Journal of Finance} 48(1), 1993.
\item W.~F.~M.~De Bondt et R.~Thaler, ``Does the stock market overreact?'', \emph{Journal of Finance} 40(3), 1985.
\item M.~Parkinson, ``The extreme value method for estimating the variance of the rate of return'', \emph{Journal of Business} 53(1),
1980~; M.~B.~Garman et M.~J.~Klass, ``On the estimation of security price volatilities from historical data'', \emph{Journal of
Business} 53(1), 1980.
\item L.~C.~G.~Rogers et S.~E.~Satchell, ``Estimating variance from high, low and closing prices'', \emph{Annals of Applied
Probability} 1(4), 1991~; D.~Yang et Q.~Zhang, ``Drift-independent volatility estimation based on high, low, open, and close prices'',
\emph{Journal of Business} 73(3), 2000.
\item Y.~Amihud, ``Illiquidity and stock returns: cross-section and time-series effects'', \emph{Journal of Financial Markets} 5(1),
2002~; S.~Gervais, R.~Kaniel et D.~H.~Mingelgrin, ``The high-volume return premium'', \emph{Journal of Finance} 56(3), 2001.
\item W.~Brock, J.~Lakonishok et B.~LeBaron, ``Simple technical trading rules and the stochastic properties of stock returns'',
\emph{Journal of Finance} 47(5), 1992.
\end{omsources}

\section{Exercices}

\begin{exercise}[$\star$]\label{exo:rs:price-and-volume-features:1}
Calculer la volatilité journalière donnée par l'\omterm{def:rs:price-and-volume-features:estimators}{estimateur de Parkinson} pour une journée où $H = 51.20$ et $L = 49.80$.
\end{exercise}

\begin{exercise}[$\star$]\label{exo:rs:price-and-volume-features:2}
Rendements mensuels de 2~\%, $-1~\%$, 4~\% et 3~\% (le plus ancien en premier). Calculer la variable de momentum 3-1 à la fin du quatrième mois,
ainsi que le rendement sur quatre mois.
\end{exercise}

\begin{exercise}[$\star$]\label{exo:rs:price-and-volume-features:3}
Une action échange USD~40 millions par jour et varie de 1.2~\% en moyenne en valeur absolue. Calculer son \omterm{def:rs:price-and-volume-features:volume}{illiquidité d'Amihud}, en pourcentage
par million de USD, ainsi que celle d'une action qui échange USD~2 millions et varie de 2.5~\%.
\end{exercise}

\begin{exercise}[$\star\star$]\label{exo:rs:price-and-volume-features:4}
Calculer les poids du croisement 3/6 selon la \cref{prop:rs:price-and-volume-features:crossover}, et leur somme.
\end{exercise}

\begin{exercise}[$\star\star$]\label{exo:rs:price-and-volume-features:5}
L'\omterm{def:rs:price-and-volume-features:estimators}{estimateur de Parkinson} est 5.2 fois plus efficace qu'un rendement journalier au carré. Déterminer combien de jours d'estimations de Parkinson donnent une estimation de la variance journalière
aussi précise que soixante jours de rendements au carré close-to-close.
\end{exercise}

\begin{exercise}[$\star\star$]\label{exo:rs:price-and-volume-features:6}
En utilisant la \cref{prop:rs:price-and-volume-features:skip}, avec des écarts-types de la partie 12-1 et du rendement du dernier mois dans le rapport
$\sqrt{11}:1$, des IC de 0.03 pour la première et $-0.03$ pour la seconde, et les deux parties non corrélées, calculer l'IC du rendement complet sur
douze mois.
\end{exercise}

\begin{exercise}[$\star\star\star$]\label{exo:rs:price-and-volume-features:7}
\emph{Programmation.} Relancer \texttt{range\_race} avec 30 pas par jour au lieu de 390. Mesurer de combien le biais de surveillance discrète de Parkinson
augmente, et en expliquer la raison.
\end{exercise}

\begin{exercise}[$\star\star\star$]\label{exo:rs:price-and-volume-features:8}
\emph{Trouver la faille.} «~Notre prédicteur de \omterm{def:rs:price-and-volume-features:volume}{surprise de volume} divise le volume du jour par le volume moyen des 21 jours centrés sur aujourd'hui, de sorte
qu'il ne soit pas biaisé par les tendances du volume.~»
\end{exercise}

\section{Problème~: la course des estimateurs par l'étendue}

\begin{problem}[{Problème du week-end --- quel estimateur de variance journalière retenir, et quand}]\label{pb:rs:price-and-volume-features:1}
La simulation du chapitre~: 4\,000 jours, 390 pas browniens d'une minute par jour, une volatilité journalière de 1~\%, graine 5~; tour à tour sans dérive ni
écart, une dérive égale à la moitié de l'écart-type journalier, et un écart nocturne portant 30~\% de la variance journalière.

\medskip\noindent\textbf{Partie I --- Les formules.}
\begin{enumerate}
\item Pourquoi la constante de Parkinson vaut-elle $4\ln2$~?
\item Quel estimateur n'a besoin que du plus haut et du plus bas~?
\item Lequel est sans biais quelle que soit la dérive, et pourquoi~?
\item Qu'apporte Yang--Zhang, et que vaut son poids $k$ pour vingt jours~?
\end{enumerate}

\medskip\noindent\textbf{Partie II --- Efficacité.}
\begin{enumerate}[resume]
\item Quelles sont les efficacités de Parkinson, Garman--Klass et Rogers--Satchell par rapport au rendement journalier close-to-close au carré~?
\item Et celle de Yang--Zhang sur vingt jours par rapport à la variance close-to-close sur vingt jours~?
\item Combien de jours de rendements close-to-close remplace un jour de Garman--Klass~?
\end{enumerate}

\medskip\noindent\textbf{Partie III --- Biais.}
\begin{enumerate}[resume]
\item Quels sont les biais des quatre estimateurs par l'étendue en l'absence de dérive ou de gap, et d'où viennent-ils~?
\item Quel effet la dérive a-t-elle sur les estimateurs close-to-close, de Parkinson et de Rogers--Satchell~?
\item Quel effet le gap a-t-il sur les trois estimateurs par l'étendue d'une seule journée, et sur celui de Yang--Zhang~?
\item Pourquoi l'estimateur close-to-close souffre-t-il de la dérive, mais pas du gap~?
\item Comment un rebond bid--ask sur le plus haut et sur le plus bas affecte-t-il les estimateurs par l'étendue~?
\end{enumerate}

\medskip\noindent\textbf{Partie IV --- Choix.}
\begin{enumerate}[resume]
\item Pour une action présentant de forts gaps overnight, quel estimateur retenir~?
\item Pour un marché ouvert 24 heures sur 24, sans gap et avec de fortes tendances intrajournalières, lequel~?
\item Pour une variable de volatilité qui doit être calculée à partir des seuls cours de clôture journaliers, laquelle choisir, et à quel prix~?
\item Que faudrait-il vérifier avant d'utiliser un estimateur par l'étendue sur des données réelles~?
\item Pourquoi les gains d'efficacité comptent-ils pour un prédicteur de volatilité~?
\item Énoncer le \emph{résultat nommé}~: les efficacités des quatre estimateurs par l'étendue et la part de variance que les estimateurs sur une seule journée ne captent pas
avec un gap de 30~\%.
\item Lesquelles des variables du chapitre utilisent une estimation de volatilité~?
\item En une phrase~: que sait l'étendue que le cours de clôture ne sait pas~?
\end{enumerate}
\end{problem}

\section{Questions d'entretien}

\begin{interviewq}[$\star$ \iqroles{chercheur, trader}]\label{iq:rs:price-and-volume-features:1}
Pourquoi les signaux de momentum ignorent-ils le mois le plus récent~?
\end{interviewq}

\begin{interviewq}[$\star\star$ \iqroles{chercheur, risque}]\label{iq:rs:price-and-volume-features:2}
Comment estimeriez-vous la volatilité journalière d'une action à partir d'un mois de barres journalières, avec le plus de précision possible~?
\end{interviewq}

\begin{interviewq}[$\star\star$ \iqroles{chercheur}]\label{iq:rs:price-and-volume-features:3}
Montrez qu'un \omterm{def:rs:price-and-volume-features:filter}{croisement de moyennes mobiles} est une somme pondérée de rendements passés. Quelles sont ces pondérations~?
\end{interviewq}

\begin{interviewq}[$\star\star$ \iqroles{développeur, ingénieur ML}]\label{iq:rs:price-and-volume-features:4}
Comment testez-vous automatiquement que le pipeline de variables n'utilise jamais de données futures~?
\end{interviewq}

\begin{interviewq}[$\star\star$ \iqroles{chercheur, trader}]\label{iq:rs:price-and-volume-features:5}
Que mesure l'\omterm{def:rs:price-and-volume-features:volume}{illiquidité d'Amihud}, et quelles en sont les faiblesses~?
\end{interviewq}

\begin{interviewq}[$\star\star\star$ \iqroles{chercheur}]\label{iq:rs:price-and-volume-features:6}
Une nouvelle variable a un IC nul sur votre simulateur et un IC de 0.02 sur des données réelles. Que concluez-vous de la variable, et du simulateur~?
\end{interviewq}
